
\documentclass[addpoints]{exam}
\usepackage{amsmath}
\usepackage{amssymb}
\usepackage{color}

% \printanswers

\newcommand{\event}{Counting Techniques}
\newcommand{\tournament}{}
\def\type{0}

\setlength\fillinlinelength{\textwidth}
\CorrectChoiceEmphasis{\color{red}\bfseries}
\SolutionEmphasis{\color{red}\bfseries}

\setlength\answerclearance{2pt}
\pagestyle{headandfoot}
\runningheadrule
\runningheader{\event}{\tournament}{\ifnum\type=1 Class Set Test \else Full Name:\hspace{1cm} \fi}
\runningfootrule
\runningfooter{}{Page \thepage\ of \numpages}{}

\begin{document}
\begin{center}
    {\huge Grade 12 Data Management \\ \event
    \ifprintanswers \space Key
    \else \space \ifnum\type<2 Test \fi \ifnum\type=1 \textbf{(CLASS SET)} \fi
          \ifnum\type=2 Answer Sheet \fi
    \fi \\ \vspace{3mm}\tournament\vspace{1mm}}
\end{center}

\vspace{.25\textheight}

\begin{center}
    First Name: \ifprintanswers
    \underline{\hspace{3.5cm}}\textbf{KEY}\underline{\hspace{5.5cm}}\\\vspace{5mm}
    \else \underline{\hspace{10cm}}\\ \vspace{5mm} \fi
    Last Name: \underline{\hspace{10cm}}\\ \vspace{5mm}
\end{center}

\noindent\textbf{\underline{Directions:}}
\begin{itemize}
    \ifnum\type=1 \item \textbf{This test is a class set.} Do not write on this paper. \fi
    \item Show all work for full marks. Leave answers in exact form unless stated otherwise.
    \item The test is designed to be completed in 75 minutes.
\end{itemize}
\begin{center}
\textit{For grading use only}\\
\multirowgradetable{1}[pages]
\end{center}

\newpage
\begin{questions}

\section*{Part A --- Multiple Choice (15 marks)}
\vspace{5mm}

\question[1] The value of $0!$ is:
\begin{choices}
    \choice $0$
    \correctchoice $1$
    \choice Undefined
    \choice $-1$
\end{choices}

\question[1] A restaurant offers 4 appetizers, 6 mains, and 3 desserts. The number of possible
three-course meals is:
\begin{choices}
    \choice $13$
    \choice $36$
    \correctchoice $72$
    \choice $216$
\end{choices}

\question[1] The number of ways to arrange 7 distinct books on a shelf is:
\begin{choices}
    \choice $7$
    \choice $49$
    \choice $128$
    \correctchoice $5040$
\end{choices}

\question[1] $P(8, 3)$ equals:
\begin{choices}
    \choice $24$
    \choice $56$
    \correctchoice $336$
    \choice $512$
\end{choices}

\question[1] The number of distinct arrangements of the letters in the word \textbf{LEVEL} is:
\begin{choices}
    \choice $120$
    \correctchoice $30$
    \choice $60$
    \choice $20$
\end{choices}

\question[1] The number of ways to seat 6 people around a circular table is:
\begin{choices}
    \choice $720$
    \choice $36$
    \correctchoice $120$
    \choice $24$
\end{choices}

\question[1] $\dbinom{10}{3}$ equals:
\begin{choices}
    \choice $30$
    \choice $60$
    \choice $90$
    \correctchoice $120$
\end{choices}

\question[1] The number of ways to choose a committee of 4 from 9 people is:
\begin{choices}
    \choice $9$
    \correctchoice $126$
    \choice $3024$
    \choice $6561$
\end{choices}

\question[1] Which row of Pascal's triangle gives the coefficients of $(a+b)^5$?
\begin{choices}
    \choice Row 4
    \correctchoice Row 5
    \choice Row 6
    \choice Row 7
\end{choices}

\question[1] The coefficient of $x^3$ in the expansion of $(1+x)^7$ is:
\begin{choices}
    \choice $21$
    \correctchoice $35$
    \choice $42$
    \choice $56$
\end{choices}

\question[1] $\dbinom{n}{r} = \dbinom{n}{n-r}$ is an example of:
\begin{choices}
    \choice The multiplication principle
    \choice Pascal's identity
    \correctchoice The symmetry property of combinations
    \choice The binomial theorem
\end{choices}

\question[1] A hockey team has 15 players. The number of ways to choose a starting line-up
of 6 (where order does not matter) is:
\begin{choices}
    \choice $P(15,6)$
    \correctchoice $\dbinom{15}{6}$
    \choice $15^6$
    \choice $6!$
\end{choices}

\question[1] How many 3-digit numbers can be formed from $\{1,2,3,4,5\}$ if repetition is
not allowed?
\begin{choices}
    \choice $15$
    \correctchoice $60$
    \choice $125$
    \choice $120$
\end{choices}

\question[1] Pascal's identity states that $\dbinom{n}{r} + \dbinom{n}{r+1} =$
\begin{choices}
    \choice $\dbinom{n+1}{r}$
    \correctchoice $\dbinom{n+1}{r+1}$
    \choice $\dbinom{n}{r+1}$
    \choice $2\dbinom{n}{r}$
\end{choices}

\question[1] The number of ways to distribute 5 distinct prizes among 3 students (each student
can receive any number of prizes) is:
\begin{choices}
    \choice $15$
    \choice $60$
    \correctchoice $243$
    \choice $120$
\end{choices}


\section*{Part B --- Short Answer (33 marks)}

\question In how many ways can 8 people line up for a photograph?
\begin{parts}
    \part[1] With no restrictions.
    \part[2] If two specific people, Alex and Beth, must stand next to each other.
    \part[2] If Alex must stand at one end and Beth must stand at the other end.
    \part[3] If Alex, Beth, and Cam must all stand together (as a block).
\end{parts}
\begin{solution}[9cm]
\textbf{(a)} $8! = \mathbf{40\,320}$

\textbf{(b)} Treat Alex \& Beth as one unit: $7!$ arrangements of the 7 units, times $2!$ for
internal arrangement of the pair:
$7! \times 2 = 5040 \times 2 = \mathbf{10\,080}$

\textbf{(c)} Alex at one end (2 choices: left or right), Beth at the other (1 remaining end),
remaining 6 in $6!$ ways:
$2 \times 6! = 2 \times 720 = \mathbf{1440}$

\textbf{(d)} Treat the block of 3 as one unit: $6!$ arrangements $\times\, 3!$ internal
arrangements:
$6! \times 3! = 720 \times 6 = \mathbf{4320}$
\end{solution}

\question The letters of the word \textbf{STATISTICS} are arranged randomly.
\begin{parts}
    \part[1] How many letters are there, and how many of each type?
    \part[3] How many distinct arrangements are there?
    \part[2] How many arrangements begin and end with the letter S?
\end{parts}
\begin{solution}[8cm]
\textbf{(a)} 10 letters: S appears 3 times, T appears 3 times, I appears 2 times,
A and C each appear once.

\textbf{(b)}
\[
\frac{10!}{3!\,3!\,2!\,1!\,1!} = \frac{3\,628\,800}{6 \cdot 6 \cdot 2} = \frac{3\,628\,800}{72}
= \mathbf{50\,400}
\]

\textbf{(c)} Fix one S at each end. Remaining 8 letters: S(1), T(3), I(2), A(1), C(1).
\[
\frac{8!}{1!\,3!\,2!\,1!\,1!} = \frac{40\,320}{12} = \mathbf{3360}
\]
\end{solution}

\question A student council has 7 boys and 5 girls.
\begin{parts}
    \part[2] In how many ways can a committee of 4 be chosen with no restrictions?
    \part[3] In how many ways can a committee of 4 be chosen with exactly 2 boys and
    2 girls?
    \part[3] In how many ways can a committee of 4 be chosen with at least 1 girl?
\end{parts}
\begin{solution}[9cm]
\textbf{(a)} $\dbinom{12}{4} = \dfrac{12!}{4!\,8!} = \mathbf{495}$

\textbf{(b)} $\dbinom{7}{2}\dbinom{5}{2} = 21 \times 10 = \mathbf{210}$

\textbf{(c)} Use complement: (at least 1 girl) $=$ (all committees) $-$ (no girls):
\[
\binom{12}{4} - \binom{7}{4} = 495 - 35 = \mathbf{460}
\]
\end{solution}

\question Use the Binomial Theorem to answer the following.
\begin{parts}
    \part[3] Expand $(2x - 3)^4$ completely.
    \part[3] Find the term containing $x^2$ in the expansion of $(3x + 2)^5$.
    \part[2] Find the middle term in the expansion of $\left(x + \dfrac{1}{x}\right)^6$.
\end{parts}
\begin{solution}[9cm]
\textbf{(a)} $(2x-3)^4 = \sum_{k=0}^{4}\binom{4}{k}(2x)^{4-k}(-3)^k$:
\begin{align*}
k=0&: (2x)^4 = 16x^4\\
k=1&: 4(2x)^3(-3) = -96x^3\\
k=2&: 6(2x)^2(9) = 216x^2\\
k=3&: 4(2x)(−27) = -216x\\
k=4&: 81
\end{align*}
$= \mathbf{16x^4 - 96x^3 + 216x^2 - 216x + 81}$

\textbf{(b)} General term: $\dbinom{5}{k}(3x)^{5-k}(2)^k$.
For $x^2$: $5-k=2 \Rightarrow k=3$.
\[
\binom{5}{3}(3x)^2(2)^3 = 10 \cdot 9x^2 \cdot 8 = \mathbf{720x^2}
\]

\textbf{(c)} $(x+x^{-1})^6$ has 7 terms; middle term is $k=3$:
\[
\binom{6}{3}x^3 \cdot x^{-3} = 20 \cdot 1 = \mathbf{20}
\]
\end{solution}

\end{questions}
\end{document}
